\chapitre{Algèbres de Lie}

\noindent
Dans tout ce chapitre, soit $\K$ désigne $\R$ ou $\C$. \\

Les algèbres de matrices $\mathcal{M}(n, \K)$ sont muni de la multiplication matricielle qui n'est en général, non-commutative. Ainsi il est intéressant de considérer le commutateur $[X, Y] = XY-YX$ qui mesure ce défaut de commutativité. Il est bilinéaire, antisymétrique et vérifie l'identité de Jacobi mais il n'est ni associatif ni compatible avec la structure multiplicative de l'algèbre. Une algèbre de Lie est précisément l'abstraction de cette structure : un espace vectoriel muni d'un crochet bilinéaire, antisymétrique et vérifiant l'identité de Jacobi, sans nécéssairement de structure d'algèbre sous-jacente. \\
 
\section{Généralités}
\subsection{Définitions et premières propriétés}

\begin{definition}
    Une \emph{algèbre de Lie} $\mathfrak{g}$ est un $\K$-espace vectoriel munie d'une application bilinéaire $\mathfrak{g}\times \mathfrak{g} \to \mathfrak{g}$, notée $(x,y) \mapsto [x, y]$ et appelée \emph{crochet de Lie}, vérifiant:
    \begin{ilist}
        \item antisymétrie\footnote{si le corps est de caractéristique différente de $2$, cet axiome est équivalent à demander que $\forall x \in \mathfrak{g}, \quad [x,x] = 0$}: $\forall x, y \in \mathfrak{g}, \quad [x, y] = - [y, x]$
        \item identité de Jacobi\index{identité de Jacobi}: $\forall x, y, z \in \mathfrak{g}, \quad [x, [y, z]] + [y, [z, x]] + [z, [x, y]] = 0$
    \end{ilist}
\end{definition}

Les principales algèbres de Lie qui nous intéresseront sont des algèbres de Lie dites \emph{matricielles} c'est à dire des sous-algèbres de Lie de $\mathcal{M}(n, \K)$ munie du commutateur, qu'on notera dorénavant $\mathfrak{gl}(n, \K)$ pour insister sur sa structure d'algèbre de Lie et qu'on appelle l'\emph{algèbre de Lie générale linéaire}. De la même façon, pour $V$ un $\K$-espace vectoriel quelconque, on notera $\mathfrak{gl}(V)$ pour $\End(V)$\footnote{Si $V$ est de dimension finie $n$, l'identification naturelle des algèbres $\End(V)$ et $\mathcal{M}(n, \K)$ après un choix de base, induit canoniquement une identification naturelle des algèbres de Lie $\mathfrak{gl}(V)$ et $\mathfrak{gl}(n, \K)$.} muni du commutateur des endomorphismes d'espaces vectoriels. \\
Ainsi on appelle 
\begin{ilist}
    \item \emph{algèbre de Lie spéciale linéaire} la sous-algèbre de Lie 
    $$\mathfrak{sl}(n, \K) = \ker(\tr) = \{ X \in \mathfrak{gl}(n, \K) \mid \tr(X) = 0 \}$$
    \item \emph{algèbre de Lie simplectique} la sous-algèbre de Lie 
    $$\mathfrak{sp}(2n, \K) = \left\{ X \in \mathfrak{gl}(2n, \K) \mid \transpose{X} J + J X = 0 \right\}$$
    où $J = \begin{pmatrix} 0 & I_n \\ -I_n & 0 \end{pmatrix}$ est la matrice symplectique standard.
    \item \emph{algèbre de Lie orthogonale} ou \emph{spéciale orthogonale} la sous-algèbre de Lie réelle
    $$\mathfrak{o}(n) = \mathfrak{so}(n) = \{ X \in \mathfrak{gl}(n, \mathbb{R}) \mid \transpose{X} + X = 0 \}$$
    Plus généralement, pour une forme de signature\footnote{La convention choisie ici place les valeurs négatives en premier dans $I_{p,q}$. L'autre convention, qui place les valeurs positives en premier, est également répandue et donne une algèbre de Lie isomorphe.} $(p, q)$, on définit
    $$\mathfrak{o}(p, q) = \mathfrak{so}(p, q) = \{ X \in \mathfrak{gl}(p+q, \mathbb{R}) \mid \transpose{X} I_{p,q} + I_{p,q} X = 0 \}$$
    où $I_{p,q} = \begin{pmatrix} -I_p & 0 \\ 0 & I_q \end{pmatrix}$.
    \item \emph{algèbre de Lie unitaire} la sous-algèbre de Lie réelle 
    $$\mathfrak{u}(n) = \{ X \in \mathfrak{gl}(n, \C) \mid X^\dagger + X = 0 \}$$
    où $X^\dagger = \transpose{{\overline{X}}}$ désigne le conjugué hermitien. 
    \item \emph{algèbre de Lie spéciale unitaire} est la sous-algèbre de Lie réelle
    $$\mathfrak{su}(n) = \{ X \in \mathfrak{u}(n) \mid \tr(X) = 0 \}$$
\end{ilist}

\begin{definition}
    Un \emph{morphisme d'algèbres de Lie} est une application linéaire $\varphi: \mathfrak{g} \to \mathfrak{h}$ telle que 
    $$\forall x, y \in \mathfrak{g}, \quad \varphi([x,y]) = [\varphi(x), \varphi(y)]$$
\end{definition}

Un exemple fondamental de morphisme d'algèbres de Lie est la \emph{représentation adjointe} : pour tout $x\in \mathfrak{g}$ on définie l'application linéaire $\ad_x : \mathfrak{g} \to \mathfrak{g}$ par 
$$\ad_x(y) = [x, y] \quad \forall y \in \mathfrak{g}$$
Ainsi on définit l'application linéaire $\ad : \mathfrak{g} \to \mathfrak{gl}(\mathfrak{g})$ qui d'après l'identité de Jacobi est un morphisme d'algèbres de Lie.

\begin{proposition}
    Soit $\varphi : \mathfrak{g} \to \mathfrak{h}$ un morphisme d'algèbre de Lie bijectif. Alors $\varphi^{-1} : \mathfrak{h} \to \mathfrak{g}$ est un morphisme d'algèbre de Lie. 
\end{proposition}

\omitted

\begin{definition}
    Une \emph{sous-algèbre de Lie} de $\mathfrak{g}$ est un sous-espace vectoriel $\mathfrak{h} \subset \mathfrak{g}$ stable par le crochet de Lie, c'est à dire :
    $$\forall x, y \in \mathfrak{h}, \quad [x,y] \in \mathfrak{h}$$ $\mathfrak{h}$ est donc une algèbre de Lie pour le crochet de Lie induit.
\end{definition}

\begin{definition}
    Une \emph{algèbre de Lie abélienne} est une algèbre de Lie $\mathfrak{g}$ dont le crochet est identiquement nul :
    $$\forall x, y \in \mathfrak{g}, \quad [x, y] = 0$$
\end{definition}

\subsection{Idéaux d'une algèbre de Lie}

\begin{definition}
    Un \emph{idéal} de $\mathfrak{g}$ est un sous-espace vectoriel $I \subset \mathfrak{g}$ absorbant pour le crochet de Lie, c'est à dire :
    $$\forall x \in \mathfrak{g}, \forall y \in I, \quad [x, y] \in I$$ 
\end{definition}

\begin{proposition}
    Soit $\varphi: \mathfrak{g} \to \mathfrak{h}$ un morphisme d'algèbres de Lie.
    \begin{ilist}
        \item L'image directe $\Img(\varphi)$ est une sous-algèbre de Lie de $\mathfrak{h}$.
        \item Le noyau $\ker(\varphi)$ est un idéal de $\mathfrak{g}$.
    \end{ilist}
\end{proposition}

\omitted

\begin{definition}
    Le \emph{centre} d'une algèbre de Lie $\mathfrak{g}$ est le sous-espace vectoriel
    $$\mathcal{Z}(\mathfrak{g}) = \{ x \in \mathfrak{g} \mid \forall y \in \mathfrak{g},\ [x, y] = 0 \}$$
\end{definition}

\begin{proposition}
    $\mathcal{Z}(\mathfrak{g})$ est un idéal abélien de $\mathfrak{g}$. C'est le plus grand idéal abélien de $\mathfrak{g}$.
\end{proposition}

\omitted

\begin{definition}
    Soit $\mathfrak{g}$ une algèbre de Lie et $\mathfrak{h}$ une sous-algèbre de Lie de $\mathfrak{g}$. Le \emph{normalisateur} de $\mathfrak{h}$ dans $\mathfrak{g}$ est la sous-algèbre de Lie 
    $$\mathcal{N}_\mathfrak{g}(\mathfrak{h}) = \{x\in \mathfrak{g} \mid [x, h] \in \mathfrak{h} \quad \forall h\in \mathfrak{h}\}$$
\end{definition}

\begin{proposition}
    Soit $\mathfrak{g}$ une algèbre de Lie et $\mathfrak{h}$ une sous-algèbre de Lie de $\mathfrak{g}$. Alors $\mathfrak{h}$ est un idéal de $\mathfrak{g}$ si et seulement si $\mathcal{N}_\mathfrak{g}(\mathfrak{h}) = \mathfrak{g}$.
\end{proposition}

\subsection{Algèbre de Lie quotient, théorèmes d'isomorphismes et extensions d'algèbres de Lie}

\begin{definition}
    Soit $I$ un idéal de $\mathfrak{g}$. L'\emph{algèbre de Lie quotient} $\mathfrak{g}/I$ est le $\K$-espace vectoriel quotient muni du crochet de Lie
    $$[x+I, y+I] = [x,y]+I$$
    où $x + I$ désigne la classe de $x$ dans $\mathfrak{g}/I$. Ce crochet est bien défini car $I$ est un idéal.
\end{definition}

\begin{theorem}
    Soit $\varphi : \mathfrak{g} \to \mathfrak{h}$ un morphisme d'algèbres de Lie et $I \subset \ker(\varphi)$ un idéal de $\mathfrak{g}$. Alors il existe un unique morphisme d'algèbres de Lie $\phi : \mathfrak{g}/I \to \mathfrak{h}$ tel que le diagramme suivant est commutatif, c'est à dire tel que $\varphi = \phi \circ \pi$ :
    \begin{diag}
        \mathfrak{g} \arrow[r, "\varphi"] \arrow[d, epi, "\pi"'] & \mathfrak{h} \\
        \mathfrak{g}/I \arrow[ur, dashed, "\phi"']
    \end{diag}
\end{theorem}

\omitted

\begin{corollary}
    Soit $\varphi : \mathfrak{g} \to \mathfrak{h}$ un morphisme d'algèbres de Lie. Alors
    $$\mathfrak{g}/\ker(\varphi) \cong \Img(\varphi)$$
\end{corollary}

\omitted

\begin{theorem}
    Soit $\mathfrak{h}$ une sous-algèbre de Lie de $\mathfrak{g}$ et $I$ un idéal de $\mathfrak{g}$. Alors $\mathfrak{h} + I$ est une sous-algèbre de Lie de $\mathfrak{g}$, $\mathfrak{h} \cap I$ est un idéal de $\mathfrak{h}$, et
    $$\mathfrak{h}/(\mathfrak{h} \cap I) \cong (\mathfrak{h} + I)/I$$
\end{theorem}

\omitted

\begin{theorem}
    Soit $I \subset J$ deux idéaux de $\mathfrak{g}$. Alors $J/I$ est un idéal de $\mathfrak{g}/I$ et
    $$(\mathfrak{g}/I)/(J/I) \cong \mathfrak{g}/J$$
\end{theorem}

\omitted

\subsection{Somme directe d'algèbres de Lie}

\begin{definition}
    La \emph{somme directe} de deux algèbres de Lie $\mathfrak{g}$ et $\mathfrak{h}$ est le $\K$-espace vectoriel $\mathfrak{g} \times \mathfrak{h}$ muni du crochet de Lie
    $$[(x_1, y_1), (x_2, y_2)] = ([x_1, x_2]_{\mathfrak{g}},\ [y_1, y_2]_{\mathfrak{h}})$$
    On la note $\mathfrak{g} \oplus \mathfrak{h}$.
\end{definition}

\begin{theorem}
    La somme directe $\mathfrak{g} \oplus \mathfrak{h}$ vérifie les propriétés universelles suivantes : pour toute algèbre de Lie $\mathfrak{k}$,
    \begin{ilist}
        \item pour tout couple de morphismes $\varphi_1 : \mathfrak{k} \to \mathfrak{g}$ et $\varphi_2 : \mathfrak{k} \to \mathfrak{h}$, il existe un unique morphisme $\varphi : \mathfrak{k} \to \mathfrak{g} \oplus \mathfrak{h}$ tel que le diagramme suivant commute
        \begin{diag}
            \mathfrak{g} & \arrow[l, "\pi_1"'] \mathfrak{g}\oplus\mathfrak{h} \arrow[r, "\pi_2"] & \mathfrak{h} \\
            & \arrow[lu, "\varphi_1"] \mathfrak{k} \arrow[ru, "\varphi_2"'] \arrow[u, dashed, "\phi"']& 
        \end{diag}
        \item pour tout couple de morphismes $\psi_1 : \mathfrak{g} \to \mathfrak{k}$ et $\psi_2 : \mathfrak{h} \to \mathfrak{k}$, il existe un unique morphisme $\psi : \mathfrak{g} \oplus \mathfrak{h} \to \mathfrak{k}$ tel que le diagramme suivant commute
        \begin{diag}
            \mathfrak{g} \arrow[r, "\iota_1"] \arrow[dr, "\psi_1"'] &  \mathfrak{g}\oplus\mathfrak{h} \arrow[d, dashed, "\psi"] & \arrow[l, "\iota_2"'] \mathfrak{h}  \arrow[dl, "\psi_2"] \\
            & \mathfrak{k} & 
        \end{diag}
    \end{ilist}
    où $\pi_1, \pi_2$ sont les projections canoniques et $\iota_1, \iota_2$ les injections canoniques.
\end{theorem}

\omitted

\begin{corollary}
    Les sous-espaces $\mathfrak{g} \times \{0\}$ et $\{0\} \times \mathfrak{h}$ sont des idéaux de $\mathfrak{g} \oplus \mathfrak{h}$, isomorphes respectivement à $\mathfrak{g}$ et $\mathfrak{h}$. Réciproquement, si une algèbre de Lie $\mathfrak{k}$ admet deux idéaux $I_1, I_2$ tels que $I_1 \cap I_2 = 0$ et $I_1 + I_2 = \mathfrak{k}$, alors $\mathfrak{k} \cong I_1 \oplus I_2$.
\end{corollary}

\omitted

\begin{proposition}
    La somme directe est associative à isomorphisme unique près : pour toutes algèbres de Lie $\mathfrak{g}$, $\mathfrak{h}$, $\mathfrak{k}$,
    $$(\mathfrak{g} \oplus \mathfrak{h}) \oplus \mathfrak{k} \cong \mathfrak{g} \oplus (\mathfrak{h} \oplus \mathfrak{k})$$
    On peut donc parler sans ambiguïté de la somme directe $\mathfrak{g}_1 \oplus \cdots \oplus \mathfrak{g}_n$ d'un nombre fini d'algèbres de Lie.
\end{proposition}

\omitted

\subsection{Algèbre de Lie des dérivations}

\begin{definition}
    Soit $A$ un $\K$-espace vectoriel muni d'une application bilinéaire $\mu : A \times A \to A$. Une \emph{dérivation} de $(A, \mu)$ est une application linéaire $\delta : A \to A$ vérifiant la règle de Leibniz :
    $$\delta(\mu(a,b)) = \mu(\delta(a), b) + \mu(a, \delta(b))$$
    L'ensemble des dérivations de $A$ est noté $\Der_\mu(A)$. \\
    Si $A$ est une $\K$-algèbre, une dérivation de $A$ vérifie $$\delta(ab) = \delta(a)b + a\delta(b)$$
    Si $\mathfrak{g}$ est une algèbre de Lie, une dérivation de $\mathfrak{g}$ vérifie 
    $$\delta([x,y]) = [\delta(x), y] + [x, \delta(y)]$$
\end{definition}

\begin{proposition}
    $\Der(A)$ est une sous-algèbre de Lie de $\mathfrak{gl}(A)$.
\end{proposition}

\omitted

\begin{proposition}
    Pour tout $x \in \mathfrak{g}$, $\ad x$ est une dérivation de $\mathfrak{g}$. Une telle dérivation est appelée \emph{dérivation intérieure}. Les dérivations intérieures forment un idéal de $\Der(\mathfrak{g})$, noté $\mathrm{Int}(\mathfrak{g})$.
\end{proposition}

\omitted

\subsection{Reduction et extension des scalaires}

\noindent
Les constructions qui suivent sont spécifiques à l'inclusion $\R \hookrightarrow \C$ et permettent de passer d'une algèbre de Lie réelle à une algèbre de Lie complexe, et réciproquement.

\begin{definition}
    Soit $\mathfrak{g}$ une algèbre de Lie complexe. La \emph{restriction des scalaires} de $\mathfrak{g}$ est l'algèbre de Lie réelle $\mathfrak{g}^\R$ obtenue en oubliant la structure de $\C$-espace vectoriel de $\mathfrak{g}$ pour ne retenir que sa structure de $\R$-espace vectoriel. Le crochet de Lie est inchangé. On a
    $$\dim_\R \mathfrak{g}^\R = 2 \dim_\C \mathfrak{g}$$
\end{definition}

\begin{definition}
    Soit $\mathfrak{g}$ une algèbre de Lie réelle. La \emph{complexification} de $\mathfrak{g}$ est l'algèbre de Lie complexe
    $$\mathfrak{g}^\C = \mathfrak{g} \otimes_\R \C$$
    munie du crochet de Lie étendu par $\C$-bilinéarité :
    $$[x \otimes z,\ y \otimes w] = [x, y] \otimes zw$$
    La multiplication scalaire complexe est donnée par $\lambda \cdot (x \otimes z) = x \otimes (\lambda z)$. Tout élément de $\mathfrak{g}^\C$ s'écrit de manière unique sous la forme $x \otimes 1 + y \otimes i$, que l'on note $x + iy$ avec $x, y \in \mathfrak{g}$.
\end{definition}

\begin{definition}
    Soit $\mathfrak{g}$ une algèbre de Lie complexe. Une \emph{forme réelle} de $\mathfrak{g}$ est une sous-algèbre de Lie réelle $\mathfrak{h} \subset \mathfrak{g}^\R$ telle que
    $$\mathfrak{h}^\C \cong \mathfrak{g}$$
\end{definition}

\section{Algèbres de Lie nilpotentes et résolubles}



\subsection{Algèbres de Lie nilpotentes}

\begin{definition}
    La \emph{série centrale descendante} d'une algèbre de Lie $\mathfrak{g}$ est la suite d'idéaux de $\mathfrak{g}$ définie par :
    $$\mathfrak{g}_0 = \mathfrak{g} \quad \text{et} \quad \mathfrak{g}_{i+1} = [\mathfrak{g}, \mathfrak{g}_i]$$
\end{definition}

\begin{definition}
    Une \emph{algèbre de Lie nilpotente} est une algèbre de Lie $\mathfrak{g}$ telle que sa série centrale descendante est nulle à partir d'un certain rang. C'est à dire qu'il existe $i \in \N$ tel que $\mathfrak{g}_i = 0$.
\end{definition}

\begin{proposition}
    Soit $\mathfrak{g}$ une algèbre de Lie nilpotente.
    \begin{ilist}
        \item Toute sous-algèbre de Lie de $\mathfrak{g}$ est nilpotente.
        \item Toute image directe de $\mathfrak{g}$ par un morphisme d'algèbres de Lie est nilpotente.
        \item Si $I$ est un idéal nilpotent de $\mathfrak{g}$ tel que $\mathfrak{g}/I$ est nilpotente, alors $\mathfrak{g}$ est nilpotente.
    \end{ilist}
\end{proposition}

\omitted

\subsection{Théorème d'Engel}

\begin{theorem}
    Une algèbre de Lie $\mathfrak{g}$ est nilpotente si et seulement si 
    $$\forall x\in \mathfrak{g}, \quad \ad x \quad \text{est nilpotent}$$ 
\end{theorem}

\omitted

\subsection{Algèbres de Lie résolubles}

\begin{definition}
    La \emph{série dérivée} d'une algèbre de Lie $\mathfrak{g}$ est la suite d'idéaux de $\mathfrak{g}$ définie par :
    $$\mathfrak{g}^0 = \mathfrak{g} \quad \text{et} \quad \mathfrak{g}^{i+1} = [\mathfrak{g}^i, \mathfrak{g}^i]$$
\end{definition}

\begin{definition}
    Une \emph{algèbre de Lie résoluble} est une algèbre de Lie $\mathfrak{g}$ telle que sa série dérivée est nulle à partir d'un certain rang. C'est à dire qu'il existe $i \in \N$ tel que $\mathfrak{g}^i = 0$.
\end{definition}

\begin{proposition}
    Soit $\mathfrak{g}$ une algèbre de Lie résoluble.
    \begin{ilist}
        \item Toute sous-algèbre de Lie de $\mathfrak{g}$ est résoluble.
        \item Toute image directe de $\mathfrak{g}$ par un morphisme d'algèbres de Lie est résoluble.
        \item Si $I$ est un idéal résoluble de $\mathfrak{g}$ tel que $\mathfrak{g}/I$ est résoluble, alors $\mathfrak{g}$ est résoluble.
        \item Si $I$ et $J$ sont deux idéaux résolubles de $\mathfrak{g}$, alors $I + J$ est résoluble.
    \end{ilist}
\end{proposition}

\omitted

\begin{proposition}
    Toute algèbre de Lie nilpotente est résoluble.
\end{proposition}

\omitted

\begin{definition}
    Le \emph{radical} d'une algèbre de Lie $\mathfrak{g}$, noté $\rad(\mathfrak{g})$, est le plus grand idéal résoluble de $\mathfrak{g}$.
\end{definition}

\subsection{Théorème de Lie}

\begin{theorem}
    Soit $\mathfrak{g}$ une sous-algèbre de Lie complexe et résoluble de $\mathfrak{gl}(V)$ où $V$ est un $\C$-espace vectoriel non nul de dimension finie. Alors $V$ possède un vecteur propre commun à tous les élements de $\mathfrak{g}$
\end{theorem}

\omitted

\subsection{Critère de Cartan}

\begin{theorem}
    Soit $\mathfrak{g}$ une sous-algèbre de $\mathfrak{gl}(V)$, où $V$ est un $\K$-espace vectoriel de dimension finie, telle que $\tr(xy)=0$ pour tout $x\in [\mathfrak{g}, \mathfrak{g}]$ et tout $y\in\mathfrak{g}$. Alors $\mathfrak{g}$ est résoluble.
\end{theorem}

\omitted

\section{Algèbres de Lie semi-simples}

\subsection{Définitions et premières propriétés}

\noindent Dans cette partie on suppose que $K$ est de caractéristique nulle.

\begin{definition}
    Une \emph{algèbre de Lie simple} est une algèbre de Lie $\mathfrak{g}$ dont le centre est trivial et qui ne possède aucun idéal autres que $0$ et elle-même.
\end{definition}

\begin{definition}
    une \emph{algèbre de Lie semi-simple} est une algèbre de Lie $\mathfrak{g}$ telle que $\rad(\mathfrak{g}) = 0$.
\end{definition}

\begin{proposition}
    Soit $\mathfrak{g}$ une algèbre de Lie. Alors $\mathfrak{g}/\rad(\mathfrak{g})$ est une algèbre de Lie semi-simple.
\end{proposition}

\omitted

\subsection{Forme de Killing}

\begin{definition}
    Soit $\mathfrak{g}$ une algèbre de Lie. La \emph{forme de Killing} de $\mathfrak{g}$ est l'application $\kappa$ définie par:
    $$\forall x, y\in \mathfrak{g}, \quad \kappa(x,y) = \tr(\ad x \ad y)$$
\end{definition}

\begin{proposition}
    Soit $\mathfrak{g}$ une algèbre de Lie. Alors $\mathfrak{g}$ est semi-simple si et seulement si sa forme de Killing est non-dégénérée.
\end{proposition}

\omitted

\begin{proposition}
    Soit $\mathfrak{g}$ une algèbre de Lie semi-simple. Alors toute dérivation de $\mathfrak{g}$ est intérieure :
    $$\Der(\mathfrak{g}) = \mathrm{Int}(\mathfrak{g})$$
\end{proposition}

\omitted

\subsection{Décomposition en somme directe d'idéaux simples}

\begin{proposition}
    Une algèbre de Lie $\mathfrak{g}$ est semi-simple si et seulement si elle se décompose en somme directe d'idéaux simples :
    $$\mathfrak{g} \cong \mathfrak{g}_1 \oplus \cdots \oplus \mathfrak{g}_n$$
    où les $\mathfrak{g}_i$ sont des idéaux simples de $\mathfrak{g}$, uniques à permutation près.
\end{proposition}

\omitted

\subsection{sous-algèbre de Cartan}

\begin{definition}
    Soit $\mathfrak{g}$ une algèbre de Lie. Une \emph{sous-algèbre de Cartan} est une sous-algèbre de Lie $\mathfrak{h}$ nilpotente et égale à son normalisateur.
\end{definition}

\begin{proposition}
    Soit $\mathfrak{g}$ une algèbre de Lie complexe semi-simple. Alors les sous-algèbre de Cartan sont exactement les sous-algèbres abéliennes maximales formées d'éléments $\ad$-diagonalisables.
\end{proposition}

\omitted

\subsection{Système de racines}

Dans cette section soit $\mathfrak{g}$ une algèbre de Lie complexe semi-simple et soit $\mathfrak{h}$ une sous-algèbre de Cartan de $\mathfrak{g}$.

\begin{proposition}
    On dispose de la décomposition suivante 
    $$\mathfrak{g} = \bigoplus_{\alpha \in \mathfrak{h}^*}\mathfrak{g}_\alpha$$
    où $\mathfrak{g}_\alpha = \{x\in \mathfrak{g} \mid \ad_x(h) = \alpha(h)x \quad \forall h \in \mathfrak{h}\}$
\end{proposition}

\omitted

\begin{definition}
    On appelle \emph{racine} tout $\alpha \in \mathfrak{h}^*$ non nulle telle que $\mathfrak{g}_\alpha \neq 0$. On note $\Phi$ l'ensemble des racines de $\mathfrak{g}$.
\end{definition}